\section{Εισαγωγή}
Η ανάλυση των ηλεκτρομηχανικών ταλαντώσεων αποτελεί βασικό εργαλείο για την αξιολόγηση της ευστάθειας των Συστημάτων Ηλεκτρικής Ενέργειας, καθώς επιτρέπει τον εντοπισμό των κυρίαρχων ταλαντωτικών ρυθμών που εμφανίζονται μετά από διαταραχές.

Στην παρούσα ενότητα εξετάζεται η εφαρμογή της μεθόδου ταυτοποίησης Matrix Pencil (MP) σε αποκρίσεις \textit{ringdown} που προκύπτουν από δυναμικές προσομοιώσεις στο περιβάλλον \textit{PowerFactory}. Η μεθοδολογία εφαρμόζεται σε χρονοσειρές τάσης ($V$), ρεύματος ($I$), ενεργού ($P$) και άεργου ($Q$) ισχύος, μετά από κατάλληλη προεπεξεργασία, με στόχο την εξαγωγή και τον χαρακτηρισμό των κυρίαρχων ταλαντωτικών ιδιοτιμών και την αξιολόγηση της ποιότητας της ανακατασκευής των σημάτων.

Η διερεύνηση πραγματοποιείται σε δύο Συστήματα Ηλεκτρικής Ενέργειας διαφορετικής κλίμακας και πολυπλοκότητας, στο \textit{Kundur Two-Area System}~\cite{kundur1994} και στο \textit{IEEE 39-Bus System}~\cite{digsilent39,pai1989}. Με αυτόν τον τρόπο, η αναφορά δεν περιορίζεται μόνο σε ένα κλασικό δοκιμαστικό σύστημα αναφοράς, αλλά επιτρέπει τη συγκριτική αξιολόγηση της συμπεριφοράς της μεθόδου ως προς την αναγνώριση \textit{inter-area} και \textit{intra-area} modes.

\section{Διαδικασία Παραγωγής Δεδομένων}
\subsection{Γενική Μεθοδολογία}
Το πρώτο βήμα περιελάμβανε την εισαγωγή των δύο ΣΗΕ στο περιβάλλον του \textit{PowerFactory} για τη δημιουργία των απαραίτητων χρονοσειρών. Συγκεκριμένα, εκτελέστηκε δυναμική προσομοίωση στο πεδίο του χρόνου (\textit{RMS Simulation}), αφού πρώτα υπολογίστηκαν οι αρχικές συνθήκες των συστημάτων (\textit{Initial Conditions}) για την επίλυση της ροής φορτίου και την εξασφάλιση της μόνιμης κατάστασης λειτουργίας. Η συνολική διάρκεια της προσομοίωσης ορίστηκε και στις δύο περιπτώσεις στα $50$ s, με χρονικό βήμα ολοκλήρωσης (\textit{step size}) $0.01$ s, επιτρέποντας την πλήρη καταγραφή της απόκρισης ελεύθερης ταλάντωσης (\textit{ringdown}) καθώς τα συστήματα τείνουν προς ένα νέο σημείο ισορροπίας. Για την παρακολούθηση των φαινομένων, εξήχθησαν σε μορφή αρχείων CSV οι μεταβλητές τάσης ($V$), ρεύματος ($I$), ενεργού ($P$) και άεργου ($Q$) ισχύος για καθεμία από τις γεννήτριες των συστημάτων.

\subsection{Διαδικασία Παραγωγής Δεδομένων για το Kundur Two-Area System}
 Ως συμβάν διαταραχής (\textit{event}) για τη διέγερση των ηλεκτρομηχανικών ταλαντώσεων, ορίστηκε η αποσύνδεση (\textit{outage}) μίας εκ των κρίσιμων γραμμών μεταφοράς που συνδέουν τις δύο περιοχές (\textit{tie-lines}) τη χρονική στιγμή $t = 0$ s. Η γραμμή παρέμεινε εκτός λειτουργίας για χρονικό διάστημα $0.1$ s, μετά το οποίο πραγματοποιήθηκε η αυτόματη επανασύνδεσή της ($t = 0.1$ s).

\subsection{Διαδικασία Παραγωγής Δεδομένων για το IEEE 39-Bus System}
Για την παραγωγή των δεδομένων προσομοιώθηκε μία αύξηση της ενεργού ισχύος $P$ του φορτίου $3$ (\textit{Load 03}) κατά $2\%$ τη χρονική στιγμή $t = 0$ s. Το συγκεκριμένο σενάριο επιλέχθηκε ως αντιπροσωπευτική και επαναλήψιμη περίπτωση διαταραχής φορτίου για τη συστηματική αξιολόγηση της μεθόδου στο \textit{IEEE 39-Bus System}~\cite{digsilent39,sfetkos2024measurement}.

\subsection{Προεπεξεργασία Σημάτων (\textit{Signal Pre-processing})}
Πριν την εφαρμογή του αλγορίθμου MP, τα πρωτογενή δεδομένα του \textit{PowerFactory} υποβλήθηκαν σε απαραίτητη ψηφιακή επεξεργασία (μέσω της γλώσσας Python), προκειμένου να εξασφαλιστεί η αξιοπιστία της ταυτοποίησης. Τα στάδια της προεπεξεργασίας περιελάμβαναν:

\begin{enumerate}
    \item \textbf{Αποκοπή Χρονοσειρών:} Διατηρήθηκαν τα δεδομένα για $t > 0.2$ s, ώστε να επικεντρωθεί η ανάλυση στο στάδιο της ελεύθερης ταλάντωσης. (\textit{ringdown}).
    \item \textbf{Αφαίρεση Γραμμικής Τάσης (\textit{Detrending}):} Αφαιρέθηκε η βέλτιστη γραμμική τάση που περιλαμβάνει τη μόνιμη συνιστώσα και τυχόν κλίση του σήματος.
    \item \textbf{Φιλτράρισμα (\textit{Filtering}):} Εφαρμόστηκε βαθυπερατό φίλτρο (\textit{Low-Pass}) με συχνότητα αποκοπής $f_{c} = 10$ Hz για την απομάκρυνση υψίσυχνου θορύβου και συνιστωσών εκτός του πεδίου ενδιαφέροντος.
\end{enumerate}

\subsection{Στάδιο ελέγχου και φιλτραρίσματος των παραγόμενων δεδομένων (Data Screening)}
\label{subsec:data_screening}

Μετά την εξαγωγή των modal estimates από τις επιμέρους μεθόδους αναγνώρισης, εφαρμόστηκε στάδιο φιλτραρίσματος (\textit{data screening}) των δεδομένων, με στόχο τη διασφάλιση της καταλληλότητάς τους για τη διαδικασία συσταδοποίησης (\textit{clustering}). Το στάδιο αυτό δεν μετέβαλε τη διαδικασία υπολογισμού των modal παραμέτρων, αλλά λειτούργησε ως βήμα προεπεξεργασίας πριν από την εφαρμογή των τεχνικών συσταδοποίησης. Η ανάγκη για data screening προέκυψε από το γεγονός ότι οι επιμέρους εκτιμήσεις που παράγονται από την ανάλυση MP ενδέχεται να περιλαμβάνουν τιμές που δεν σχετίζονται με πραγματικούς ηλεκτρομηχανικούς ρυθμούς ταλάντωσης των συστημάτων, αλλά οφείλονται σε θόρυβο, αριθμητικές αβεβαιότητες ή μη φυσικούς ρυθμούς. Για τον λόγο αυτό, πριν από την ομαδοποίηση, διατηρήθηκαν μόνο οι ιδιοτιμές με συχνότητες εντός του διαστήματος ενδιαφέροντος, το οποίο καθορίστηκε από τη φυσική συμπεριφορά των συστημάτων και τις σχετικές μελέτες της βιβλιογραφίας, ενώ απορρίφθηκαν και οι εκτιμήσεις με μη αρνητική ή πρακτικά μηδενική απόσβεση μέσω του κριτηρίου $\sigma \leq -10^{-3}$.

Οι ταλαντωτικοί ρυθμοί των ΣΗΕ διακρίνονται γενικά σε \textit{inter-area} και \textit{intra-area} modes. Οι \textit{inter-area} ταλαντώσεις αντιστοιχούν σε σχετική ταλάντωση ομάδων γεννητριών που ανήκουν σε διαφορετικές περιοχές του συστήματος, ενώ οι \textit{intra-area} ταλαντώσεις αφορούν γεννήτριες που ταλαντώνονται εντός της ίδιας περιοχής. Για το \textit{Kundur Two-Area System}, σύμφωνα με τη σχετική βιβλιογραφία, οι \textit{inter-area} modes εμφανίζονται στην περιοχή από $0.1$~Hz έως $0.7$~Hz, ενώ οι \textit{intra-area} modes εντοπίζονται στην περιοχή από $0.7$~Hz έως $2.0$~Hz~\cite{estimation2021}. Για το \textit{IEEE 39-Bus System}, τα modes αναφοράς που αξιοποιήθηκαν στην παρούσα διερεύνηση εντοπίζονται επίσης εντός της περιοχής $0.1-2.0$~Hz~\cite{sarkar2024drga}. Συνεπώς, στο στάδιο data screening διατηρήθηκαν μόνο οι εκτιμήσεις με συχνότητα εντός του διαστήματος $[0.1,\,2.0]$~Hz και με αρνητικό ρυθμό απόσβεσης που να ικανοποιεί το όριο $\sigma \leq -10^{-3}$. Έτσι προέκυψε ένα καθαρό σύνολο δεδομένων, κατάλληλο για την εφαρμογή των τεχνικών συσταδοποίησης που αξιοποιήθηκαν στη συνέχεια της ανάλυσης.

Σημειώνεται ότι το στάδιο data screening εφαρμόστηκε αποκλειστικά πριν από τη διαδικασία clustering, δηλαδή στο πλαίσιο της εξαγωγής αντιπροσωπευτικών ομάδων ιδιοτιμών. Επομένως, το screening αποτελεί στοχευμένο στάδιο προεπεξεργασίας για τις ανάγκες της ομαδοποίησης και δεν χρησιμοποιήθηκε ως γενική τροποποίηση όλων των παραγόμενων modal estimates.



